% Figure peps_nested_cylinder_orthogonalization_projected_image.
% Owner: chapter/ch24_peps_nested_cylinder_orthogonalization.tex (unlabelled remark).
% Used by: gapped-2d-notes PP Sec. 5 step (h).
% Source: arXiv OpenAI2026PolynomialPEPS, 03-patches.tex:575--580 (the projected image
% of S_j, not its intersection with the complement of W_j); no source figure.
% Formula represented: P_W(e_1+e_2)=e_1 and (I-P_W)(e_1+e_2)=e_2;
% both subsequent projections are zero, so the square commutes.
% Ink-to-map correspondence: each node is the labelled subspace;
% each arrow is the restriction of its labelled orthogonal projection.
% This is a commutative diagram, not a tensor network. There are no
% contracted or open tensor legs; no tensor boundary signature is asserted.
\begin{center}
    \begin{tikzcd}[column sep=huge,row sep=large]
        W=\spn\{e_1\} \arrow[d,"I-P_W"'] &
        S=\spn\{e_1+e_2\} \arrow[l,"P_W"'] \arrow[d,"I-P_W"] \\
        \{0\} & (I-P_W)S=\spn\{e_2\} \arrow[l,"P_W"]
    \end{tikzcd}
\end{center}
